\documentclass[a4paper,12pt]{article}
\usepackage{geometry}
\usepackage{longtable}
\usepackage{hyperref}
\usepackage{titlesec}
\usepackage{booktabs}

\geometry{margin=1in}
\titleformat{\section}{\Large\bfseries}{\thesection}{1em}{}

\title{\textbf{Summary of Second and Third Generation Nuclear Reactors}}
\author{Generated by ChatGPT}
\date{\today}

\begin{document}

\maketitle

\section{Introduction}
This document provides a structured summary of the main types and characteristics of \textbf{Generation II} and \textbf{Generation III} nuclear reactors, including their design categories, representative models, and notable operating countries. These generations represent the evolution of commercial nuclear power technology from the 1960s to the present day.

\section{Generation II Reactors}

\subsection{Overview}
Generation II reactors were developed primarily between the 1960s and 1990s. They improved upon first-generation prototypes with higher reliability, scalability, and safety. Most of the world’s operating reactors today are of this generation.

\subsection{Main Reactor Types}

\begin{longtable}{p{3cm}p{3cm}p{3cm}p{4cm}p{3cm}}
\toprule
\textbf{Type} & \textbf{Coolant} & \textbf{Moderator} & \textbf{Example Designs} & \textbf{Countries} \\
\midrule
PWR & Light Water & Light Water & Westinghouse PWR, VVER-1000 & USA, France, Russia \\
BWR & Light Water & Light Water & GE BWR/3--6 & USA, Japan, Sweden \\
HWR / PHWR & Heavy Water & Heavy Water & CANDU, Indian PHWR & Canada, India \\
AGR & CO$_2$ Gas & Graphite & AGR Series & UK \\
FBR & Liquid Sodium & None & BN-600, Phénix & Russia, France \\
RBMK & Light Water & Graphite & RBMK-1000 & Russia, Lithuania, Ukraine \\
\bottomrule
\end{longtable}

\subsection{Notable Features}
\begin{itemize}
    \item Typical design life: 40 years
    \item Active safety systems
    \item Operated with enriched uranium or natural uranium (for HWRs)
    \item Foundation for most modern reactor designs
\end{itemize}

\section{Generation III Reactors}

\subsection{Overview}
Generation III and III+ reactors, developed since the 1990s, include advanced safety features (often passive), improved fuel efficiency, and longer operational lifetimes (60+ years). These designs represent evolutionary improvements over Generation II reactors.

\subsection{Main Reactor Designs}

\subsubsection{Pressurized Water Reactors (PWRs)}
\begin{itemize}
    \item \textbf{AP1000} (Westinghouse, USA): Passive safety, $\sim$1100 MWe
    \item \textbf{EPR / European Pressurized Reactor} (Framatome, France): $\sim$1600 MWe, double containment
    \item \textbf{APR-1400} (KHNP, South Korea): Derived from System 80+, $\sim$1400 MWe
    \item \textbf{VVER-1200 / AES-2006} (Rosatom, Russia): Passive safety, 1200 MWe
    \item \textbf{HPR1000 / Hualong One} (China): Based on French and Chinese designs, $\sim$1100 MWe
    \item \textbf{Atmea1} (MHI/Framatome): $\sim$1200 MWe, mid-size
    \item \textbf{APWR} (Japan): $\sim$1500 MWe, high efficiency
\end{itemize}

\subsubsection{Boiling Water Reactors (BWRs)}
\begin{itemize}
    \item \textbf{ABWR} (GE-Hitachi): First Gen III reactor, digital control systems
    \item \textbf{ESBWR} (GE-Hitachi): Fully passive safety, $\sim$1520 MWe
\end{itemize}

\subsubsection{Heavy Water Reactors (HWRs)}
\begin{itemize}
    \item \textbf{ACR-1000} (Canada): Hybrid coolant system, not deployed
    \item \textbf{IPHWR-700} (India): Enhanced safety, 700 MWe, under construction
\end{itemize}

\subsubsection{Fast Neutron Reactors (FNRs)}
\begin{itemize}
    \item \textbf{BN-800} (Russia): Sodium-cooled fast reactor, operating
    \item \textbf{CEFR} (China): Experimental fast reactor
\end{itemize}

\subsection{Additional and Regional Designs}
\begin{longtable}{p{4cm}p{3cm}p{3cm}p{5cm}}
\toprule
\textbf{Reactor Design} & \textbf{Type} & \textbf{Country} & \textbf{Notes} \\
\midrule
KERENA & BWR & Germany & Passive safety concept \\
APR+ & PWR & South Korea & Upgraded APR-1400 \\
VVER-TOI & PWR & Russia & Advanced VVER-1200 variant \\
ACP100 (Linglong One) & SMR / PWR & China & Gen III+ small modular reactor \\
SMART & SMR / PWR & South Korea & 330 MWt design \\
ACR-700 & HWR & Canada & Scaled-down version of ACR-1000 \\
\bottomrule
\end{longtable}

\subsection{Key Generation III Features}
\begin{itemize}
    \item Passive and active redundant safety systems
    \item 60+ year design lifetime (extendable to 80 years)
    \item 10--15\% higher thermal efficiency than Gen II
    \item Extended fuel burnup (45--60 GWd/tU)
    \item Double containment and improved seismic tolerance
\end{itemize}

\section{Conclusion}
Generation II reactors laid the foundation for modern nuclear energy, emphasizing reliability and scalability. Generation III and III+ designs build upon these systems with a focus on safety, longevity, and efficiency, representing the most advanced large-scale commercial nuclear reactors in operation today.

\end{document}
